\documentclass[10pt,a4paper]{amsart}
\usepackage[margin=19mm]{geometry}
\usepackage{amsmath,amssymb,mathtools}
\usepackage[scr=rsfs,cal=euler]{mathalfa}
\usepackage{xfrac}
\usepackage[unicode,hidelinks,bookmarks=false]{hyperref}
\newcommand{\ie}{i.e.\ }
\newcommand{\cf}{cf.\ }
\newcommand{\resp}{respectively\ }
\newcommand{\Subsection}{\subsection}
\providecommand{\nobreakdash}{\nobreak\hbox{-}}
\setlength{\emergencystretch}{2em}
\multlinegap0pt
\usepackage{bm}
\usepackage{bbm}
\usepackage{dsfont}
\usepackage{xspace}
\usepackage{cancel}
\usepackage{mathtools}
\usepackage{amsthm}
\makeatletter
\@ifundefined{theorem}{\newtheorem{theorem}{Theorem}}{}
\@ifundefined{lemma}{\newtheorem{lemma}[theorem]{Lemma}}{}
\makeatother
\renewcommand\ge{\geq}
\renewcommand\leq{\leqslant}
\title[Relative differential closure in Hardy fields]{Relative differential closure in Hardy fields}
\author{Matthias Aschenbrenner, Lou van den Dries, Joris van der Hoeven}
\date{Source: arXiv:2412.10764v2 (CC BY 4.0)}
\begin{document}
\maketitle

\noindent\emph{Source and licence.} Relative differential closure in Hardy fields. Version arXiv:2412.10764v2 (CC BY 4.0), distributed under the Creative Commons Attribution 4.0 International licence, as recorded in the arXiv licence metadata for this exact version and shown on the abstract page https://arxiv.org/abs/2412.10764v2. Full rights notice: licenses/arxiv-2412.10764-CC-BY-4.0.txt.\par\smallskip

\noindent\textit{Editorial scope.} Counted unit: the Lemma whose source label is \texttt{None} in the article (source0.tex lines 1062--1066, source label \texttt{None}), delivered together with the article's own complete original proof of it (source0.tex lines 1067--1088, one \texttt{proof} environment of 22 lines). No mathematics is written, repaired or re-derived: statement and proof are the article's own text line for line. Required context retained uncounted: none. Every definition, display and label of those lines is retained unchanged. Omitted from this package and not redistributed: 1--1061, 1089--2080. Nothing inside a retained excerpt is omitted or altered: every retained body is delivered exactly as the article's lines are written, so no omission receipt of any kind accompanies this package. Citations: the retained excerpts carry no citation command at all, so no bibliography entry of the article is transcribed and no reference is redistributed. Reference adaptation: none was needed and none was made; every reference inside the delivered unit resolves inside it, so no unresolved reference or citation is produced. Printed slips kept literal: no printed word, symbol, hypothesis or number of the counted statement or of its proof was altered. Layout and class: the article's own class is \texttt{amsart} with packages including upgreek, graphicx, enotez, amssymb, epic, mathrsfs, accents, array, stmaryrd, enumitem, longtable, hyperref, etoolbox, tikz; the wrapper is the lane's pinned \texttt{amsart} editorial wrapper at 10pt on A4, which restates the theorem-like environments the retained text needs and adds the article's own macro definitions that the retained text needs (\texttt{\textbackslash ge}, \texttt{\textbackslash leq}), with extra wrapper packages (bm, bbm, dsfont, xspace, cancel, mathtools). The delivered numbering is the unit's own. Assets: no retained span contains a figure environment, a graphics-inclusion command, a PDF-page inclusion, a table or a TikZ picture; no image file is copied into this package, the package declares no asset dependency and no image file is redistributed with it. Licence: version arXiv:2412.10764v2 of this article is distributed under the Creative Commons Attribution 4.0 International licence, as recorded in the arXiv licence metadata for this exact version and shown on the abstract page https://arxiv.org/abs/2412.10764v2. A full rights notice is delivered at licenses/arxiv-2412.10764-CC-BY-4.0.txt. Characters: every retained excerpt is pure ASCII, so the delivered engine, XeLaTeX with its default Latin Modern text font, reports no missing character.
\begin{lemma}
Suppose $K$ is not algebraically closed, $C\ne K$, and
$K$ is weakly $r$-$\d$-closed in $L$. Let $Q_1,\dots,Q_m\in K\{Y\}^{\neq}$ of order~$\leq r$ have a common zero
in~$L$,~$m\ge 1$. Then they have a common zero in $K$.
\end{lemma}

\begin{proof}
We claim that some polynomial $\Phi\in K[X_1,\dots,X_m]$ has $(0,\dots,0)\in K^m$ as its only
zero in $K^m$.
(This is folklore, but we didn't find a proof in the literature.)
The claim is clear for $m=1$. Since $K$ is not algebraically closed, we have a univariate polynomial $f\in K[X]$ of degree $>1$ 
%$$f(X)\ =\  X^d + a_1 X^{d-1} + \cdots + a_d \qquad d>1,\  a_1,\dots, a_{d}\in K,$$ 
 without a zero in $K$. %, so $a_d \ne 0$.
Then the homogenization  $g(X,Y)\in K[X,Y]$
%$$g(X,Y)\ :=\  X^d + a_{1} X^{d-1} Y + \cdots + a_d Y^d \in K[X,Y]$$
 of~$f(X)$ has $(0,0)$ as its only zero in $K^2$. 
%be the of $f$, so $g(X,1)=f(X)$. Consider $(x,y) \in K^2$ with
%$g(x,y)=0$. If $y \ne 0$, then $f(x) = g(x,y) / y^d = 0$, which is
%impossible. If $y=0$, then $g(x,y) = x^d = 0$, whence $x = 0$.
This proves the claim for $m=2$. Now use induction on $m$: if $\Phi$ has the above property, then
 $g\big(\Phi(X_1, \ldots, X_{m}), X_{m+1}\big)$ has~$(0,\dots, 0,0)\in K^{m+1}$ as its only zero in $K^{m+1}$.


By the claim, the differential
polynomial~$P:=\Phi(Q_1,\dots,Q_m)\in K\{Y\}$ is nonzero (use [ADH, 4.2.1]) and
has order~$\leq r$. For $y\in L$ we have $$Q_1(y)=\cdots=Q_m(y)=0\quad \Longrightarrow\quad P(y)=0,$$
and for $y\in K$ the converse of this implication also holds. %we have  \[P(y)=0\quad \Longrightarrow\quad Q_1(y)=\cdots=Q_m(y)=0.\qedhere\]
\end{proof}

\end{document}
